% Anexo: herramientas de línea de comandos (ajuste headless).
% Se incluye con \input desde main.tex bajo un \chapter, así que empieza en \section.

Los ajustes también pueden hacerse \emph{sin interfaz}, desde la línea de
comandos, con el \textbf{mismo motor numérico} que la GUI
(\code{core.fit\_engine} y \code{mossbauer\_distribution}): los resultados son
idénticos a los que se obtienen de forma interactiva. Dos programas cubren los
dos tipos de ajuste:

\begin{itemize}[nosep]
  \item \file{mossbauer\_fit\_cli.py} --- ajustes \textbf{discretos}
        (singlete / doblete / sextete / modelos de relajación) a partir de una
        plantilla de sesión, con errores por bootstrap Monte Carlo, intervalos
        por verosimilitud perfilada y series con arranque en caliente.
  \item \file{fit\_bhf\_distribution\_cli.py} --- ajustes de
        \textbf{distribuciones hiperfinas}: $P(\BHF)$, $P(\dEQ)$ o
        $P(\mathrm{IS})$, todas las formas 1D (histograma, gaussiana, VBF,
        binomial, Czjzek), todos los regularizadores (Tikhonov, TV, MaxEnt), la física
        del kernel del panel de la GUI (Hamiltoniano completo, fuente
        polarizada, textura), kernels Voigt, correlaciones con el campo y la
        distribución 2D $P(\BHF,\dEQ)$, como histograma o como gaussianas
        bivariantes.
\end{itemize}

Ninguno de los dos necesita pantalla: funcionan en un clúster, dentro de una
cola de trabajos o llamados desde otro programa. Ambos devuelven el código de
salida $0$ si todo va bien, lo que facilita encadenarlos en scripts de shell.
\code{-{}-help} enumera todas las opciones con su valor por defecto.

\section{Ajustes discretos: \file{mossbauer\_fit\_cli.py}}
\label{sec:cli-discrete}

El flujo es el mismo que en la GUI: se prepara el modelo una vez de forma
interactiva, se guarda como sesión (\menu{Archivo > Guardar sesión…}) y ese
fichero JSON sirve de \textbf{plantilla} para cualquier número de espectros:

\begin{lstlisting}[language=bash]
python mossbauer_fit_cli.py \
    --template plantilla_alphaFe.json \
    --spectrum muestra.adt \
    --out muestra.fit.json
\end{lstlisting}

La plantilla contiene la definición completa del modelo: tipos de componente y
valores iniciales, parámetros fijos o libres, restricciones, perfil de línea
(Lorentz/Voigt), verosimilitud (gaussiana/Poisson), pérdida robusta, modelo de
absorbente, forma de onda del transductor (triangular/seno), número de arranques
múltiples y las opciones de calibración. El centro de doblado \emph{no} se toma
de la plantilla: se detecta en cada espectro, igual que hace la GUI. El JSON de
salida tiene la misma estructura que una sesión de la GUI (se puede abrir con
\menu{Archivo > Cargar sesión…}) más un bloque \code{batch\_fit\_result} con los
valores ajustados, sus errores y los estadísticos ($\chi^2$, $\chi^2$ reducido,
AIC, BIC).

\subsection{Opciones}

\begin{center}
\begin{tabular}{@{}llp{7.6cm}@{}}
\toprule
Opción & Argumento & Significado \\
\midrule
\code{-{}-template} & \file{fichero.json} & Plantilla de sesión (obligatoria). \\
\code{-{}-spectrum} & uno o más ficheros & Espectro(s) \file{.ws5}/\file{.adt}
  que ajustar. Con varios, la serie se ajusta con arranque en caliente
  secuencial (ver más abajo). \\
\code{-{}-out} & \file{fichero.json} & Fichero de salida (obligatorio). \\
\code{-{}-vmax} & mm/s & Sustituye el $v_{\max}$ de la plantilla (útil si la
  calibración cambia entre series). \\
\code{-{}-bootstrap} & $N$ & Errores por bootstrap Monte Carlo con $N$
  réplicas (solo con un espectro). Añade una sección \code{bootstrap} con la
  $\sigma_{\mathrm{MC}}$ de cada parámetro. \\
\code{-{}-profile-likelihood} & --- & Intervalos asimétricos $1\sigma/2\sigma$
  por verosimilitud perfilada (cruces de $\Delta\chi^2=1$ y $\Delta\chi^2=4$;
  ver el anexo matemático). Solo con un espectro. Añade una sección
  \code{profile\_likelihood}. \\
\code{-{}-export-job} & \file{fichero.JOB} & Escribe además el modelo ajustado
  como fichero \file{.JOB} de NORMOS-SITE, para compararlo en NORMOS o
  compartir la configuración. No ejecuta NORMOS. \\
\code{-{}-src-frac} & $F$ & Fracción resonante de la fuente (FSO de
  NORMOS-SITE), solo con el modelo de absorbente \code{transmission}: la
  transmisión satura en $1-F$ en lugar de en $0$. Por defecto $1$. \\
\code{-{}-line-asym} & $A$ & Asimetría de línea por interferencia (AKS de
  NORMOS-SITE): el perfil se multiplica por $1 - A\,(v-v_0)/\Gamma$. Por
  defecto $0$. \\
\code{-{}-quiet} & --- & No imprime el resumen por la salida estándar. \\
\bottomrule
\end{tabular}
\end{center}

\subsection{Convenios de NORMOS}

Tres opciones cambian a los convenios de NORMOS, para reproducir sus ajustes. El
primer valor de cada una es el de Fitbauer, que es el defecto; son las mismas
opciones que el submenú \menu{Ajuste > Opciones avanzadas de ajuste > Convenios
(compatibilidad NORMOS)} de la GUI.

\begin{center}
\begin{tabular}{@{}>{\raggedright\arraybackslash}p{4.4cm}>{\raggedright\arraybackslash}p{2.7cm}>{\raggedright\arraybackslash}p{8.0cm}@{}}
\toprule
Opción & Valores & Significado \\
\midrule
\code{-{}-sextet-pattern} & \code{alpha\_fe} $|$ \code{normos} & Posiciones de
  las líneas del sextete: el patrón publicado del $\alpha$-Fe, o derivadas de
  los momentos nucleares, como hace NORMOS-SITE (el campo difiere en unos
  $0.01$~T a $33$~T). Afecta al sextete de primer orden, al Hamiltoniano
  completo y al kernel de distribución. \\
\code{-{}-intensity-convention} & \code{depth} $|$ \code{area} & Si las razones
  \code{int1}/\code{int2} entre líneas son de profundidad o de área (los
  D13/D23 de NORMOS). Solo difieren si las anchuras relativas no son $1$. \\
\code{-{}-relax-convention} & \code{blume} $|$ \code{normos} & Relajación de dos
  estados con poblaciones desiguales: el modelo de Blume, con balance
  detallado, o la regla que aplica NORMOS-SITE con BSAT/BH0. Solo difieren con
  polarización $P > 0$. \\
\bottomrule
\end{tabular}
\end{center}

\subsection{Errores avanzados}

Con \code{-{}-bootstrap} $N$ el ajuste base se repite sobre $N$ espectros
sintéticos (con cuentas de Poisson reales cuando se conoce el factor de
normalización) y, junto al error de la covarianza, se da la desviación típica
muestral de cada parámetro libre:

\begin{lstlisting}[language=bash]
python mossbauer_fit_cli.py --template T.json --spectrum a.adt \
    --out a.fit.json --bootstrap 100
#  s1_bhf = 33.0446 +- 0.005988  [sigma_MC 0.006113]
\end{lstlisting}

Con \code{-{}-profile-likelihood} cada parámetro libre se recorre alrededor de
su óptimo reajustando los demás, y los cruces de $\Delta\chi^2=1$ y
$\Delta\chi^2=4$ dan intervalos asimétricos, válidos aunque falle la
aproximación parabólica de la covarianza:

\begin{lstlisting}[language=bash]
#  s1_bhf = 33.0446 +- 0.005988  [1-sigma profile -0.00607/+0.006066]
\end{lstlisting}

\subsection{Series con arranque en caliente}

Si se pasan varios espectros (una serie en temperatura o en composición), cada
ajuste parte de los valores convergidos del anterior: es el mismo arranque en
caliente secuencial que la herramienta de lotes de la GUI. La salida es un único
JSON con una fila por fichero (\code{status}, \code{values}, \code{errors},
\code{stats}); un fallo individual no detiene la serie:

\begin{lstlisting}[language=bash]
python mossbauer_fit_cli.py --template T.json \
    --spectrum serie_010K.adt serie_050K.adt serie_100K.adt \
    --out serie.batch.json
\end{lstlisting}

Para ajustes totalmente independientes (sin arranque en caliente) basta con
llamar al programa una vez por fichero desde la shell.

\section{Ajustes de distribución: \file{fit\_bhf\_distribution\_cli.py}}
\label{sec:cli-dist}

Este programa lee un fichero \file{.ws5}/\file{.adt} sin doblar, lo dobla (al
estilo de NORMOS, aprovechando los ficheros \file{.RES}/\file{.JOB} que lo
acompañen) y ajusta una distribución hiperfina con el mismo motor que el panel
de la GUI del capítulo~\ref{ch:distribuciones}. A diferencia de la GUI, los
parámetros del kernel $\delta$, $\dEQ$ y $\Gamma$ son \emph{fijos}: se dan con
\code{-{}-delta}/\code{-{}-quad}/\code{-{}-gamma} o se leen de esos ficheros.

\begin{lstlisting}[language=bash]
# Histograma P(BHF) clásico (Tikhonov):
python fit_bhf_distribution_cli.py muestra.adt --alpha 0.01 --nbins 50 --plot

# P(DEQ) con dobletes; P(IS) con singletes:
python fit_bhf_distribution_cli.py muestra.adt --variable quad --bmax 3
python fit_bhf_distribution_cli.py muestra.adt --variable delta \
    --bmin -1 --bmax 2

# VBF con dos gaussianas; histograma MaxEnt:
python fit_bhf_distribution_cli.py muestra.adt --shape vbf --vbf-components 2
python fit_bhf_distribution_cli.py muestra.adt --shape vbf --vbf-sigma-min 0.5 \
    --vbf-sigma-max 3
python fit_bhf_distribution_cli.py muestra.adt --reg-mode maxent

# Binomial al estilo NORMOS (DISTRI=3), 5 % de soluto, 12 vecinos:
python fit_bhf_distribution_cli.py muestra.adt --shape binomial \
    --binomial-conc 5 --binomial-start 33 --binomial-step -2.5

# P(dEQ) de Czjzek extendida (vidrio con Fe3+):
python fit_bhf_distribution_cli.py muestra.adt --variable quad --bmax 3 \
    --shape czjzek-ext

# Distribución 2D: histograma regularizado, o dos gaussianas bivariantes:
python fit_bhf_distribution_cli.py muestra.adt --dist-2d --qmin -1 --qmax 1
python fit_bhf_distribution_cli.py muestra.adt --dist-2d --shape vbf \
    --vbf-components 2 --qmin -1 --qmax 1

# P(IS, DEQ) con dobletes (vidrios, amorfos):
python fit_bhf_distribution_cli.py muestra.adt --dist-2d --pair delta-quad
\end{lstlisting}

\subsection{Distribución y forma}

\begin{center}
\begin{tabular}{@{}>{\raggedright\arraybackslash}p{4.4cm}>{\raggedright\arraybackslash}p{2.7cm}>{\raggedright\arraybackslash}p{8.0cm}@{}}
\toprule
Opción & Valores & Significado \\
\midrule
\code{-{}-variable} & \code{bhf} $|$ \code{quad} $|$ \code{delta} &
  Variable distribuida: $P(\BHF)$ (por defecto), $P(\dEQ)$ o
  $P(\mathrm{IS})$. Con \code{quad} y \code{delta} el campo queda fijo en
  \code{-{}-fixed-bhf} ($0$~T $\Rightarrow$ dobletes, o singletes si además
  \code{-{}-quad 0}). \\
\code{-{}-shape} & \code{histograma} $|$ \code{gaussiana} $|$ \code{vbf} $|$
  \code{binomial} $|$ \code{czjzek} $|$ \code{czjzek-ext} & Forma de la distribución
  (capítulo~\ref{ch:distribuciones}). \code{gaussiana} es el VBF con $N=1$ y
  línea lorentziana; \code{czjzek} y \code{czjzek-ext} solo con
  \code{-{}-variable quad} y dobletes (§\ref{sec:forma-czjzek}). \\
\code{-{}-czjzek-sigma}, \code{-{}-czjzek-delta0} & mm/s & Czjzek: $\sigma$ y
  $\Delta E_{Q,0}$ iniciales; con \code{-{}-fix-czjzek-sigma} /
  \code{-{}-fix-czjzek-delta0} quedan fijos. \\
\code{-{}-vbf-components} & $N$ & Número de gaussianas del VBF (1D y 2D). \\
\code{-{}-vbf-sigma-min}, \code{-{}-vbf-sigma-max} & $\sigma$ & Cotas de la
  anchura $\sigma$, común a todas las gaussianas, en unidades de la variable
  (eje $X$ en 2D). Solo con \code{gaussiana} y \code{vbf}; $0$ = automático
  (por defecto). \\
\code{-{}-vbf-sigma-min-y}, \code{-{}-vbf-sigma-max-y} & $\sigma$ & Lo mismo
  para el eje $Y$ en 2D (mm/s). \\
\code{-{}-reg-mode} & \code{tikhonov} $|$ \code{tv} $|$ \code{maxent} &
  Regularizador del histograma. \\
\code{-{}-alpha} & $\alpha$ & Intensidad de la regularización del
  histograma. \\
\code{-{}-edge-anchor} & $B$ & Lleva $P$ hacia $0$ en el primer y el último
  bin (BETA1/BETA2 de NORMOS-DIST), para que el histograma no acumule peso
  espurio en los bordes. $0$ = sin anclaje. \\
\code{-{}-binomial-conc} & \% & Binomial: concentración de soluto (CONC de
  NORMOS). Fija salvo con \code{-{}-fit-conc}. \\
\code{-{}-binomial-neighbors} & $n$ & Binomial: número de posiciones vecinas
  (NORMOS usa 12). \\
\code{-{}-binomial-start}, \code{-{}-binomial-step} & valores & Binomial: valor
  de la variable sin vecinos de soluto y cambio por vecino (BHF/DTB, QUA/DTQ,
  ISO/DTI de NORMOS). \\
\code{-{}-fit-conc} & --- & Binomial: ajusta también la concentración. \\
\bottomrule
\end{tabular}
\end{center}

\subsection{Kernel}

\begin{center}
\begin{tabular}{@{}>{\raggedright\arraybackslash}p{4.4cm}>{\raggedright\arraybackslash}p{2.7cm}>{\raggedright\arraybackslash}p{8.0cm}@{}}
\toprule
Opción & Valores & Significado \\
\midrule
\code{-{}-delta}, \code{-{}-quad}, \code{-{}-gamma} & mm/s & Parámetros fijos
  del kernel; por defecto se leen de los ficheros de NORMOS si existen
  ($\Gamma = 0.36$~mm/s si no). \\
\code{-{}-fixed-bhf} & T & Campo del kernel con \code{-{}-variable quad} o
  \code{delta}. \\
\code{-{}-profile} & \code{Lorentziana} $|$ \code{Voigt} & Perfil de línea del
  kernel; \code{-{}-voigt-sigma} da la anchura gaussiana instrumental
  (mm/s). \\
\code{-{}-delta-slope}, \code{-{}-quad-slope} & pendiente & Correlaciones
  lineales $\delta(H)$ y $\dEQ(H)$ del anexo~\ref{ap:mates}; $0$ = kernel
  clásico. \\
\code{-{}-d13}, \code{-{}-d23} & razones & Razones de área de las líneas
  (1,6)/(3,4) y (2,5)/(3,4) del kernel de sextetes (convenio de NORMOS; $3$ y
  $2$ = polvo). Solo histograma. \\
\code{-{}-kernel-treatment} & \code{1er\_orden} $|$ \code{hamiltoniano} &
  Sextetes de primer orden, o el promedio de polvo del Hamiltoniano completo
  (entonces \code{-{}-quad} es el módulo de un EFG aleatorio, con asimetría
  \code{-{}-kernel-eta}), el análogo exacto de EXACT/QUP en NORMOS-DIST. Solo
  histograma. \\
\code{-{}-source-polarized} & --- & Fuente polarizada (METHOD=4/POLAR de
  NORMOS-DIST): kernel en peine de 36 líneas, con \code{-{}-source-bhf},
  \code{-{}-source-theta} y \code{-{}-absorber-theta}. Solo histograma con
  \code{-{}-variable bhf}. \\
\code{-{}-sharp-bhf} & T (repetible) & Añade sextetes cristalinos (nítidos)
  sobre la distribución, p.~ej.\ \code{-{}-sharp-bhf 33 -{}-sharp-bhf 49};
  \code{-{}-sharp-delta}, \code{-{}-sharp-quad} y \code{-{}-sharp-gamma} dan
  sus parámetros (por defecto, los del kernel). \\
\bottomrule
\end{tabular}
\end{center}

\subsection{Malla, fondo y doblado}

\begin{center}
\begin{tabular}{@{}>{\raggedright\arraybackslash}p{4.4cm}>{\raggedright\arraybackslash}p{2.7cm}>{\raggedright\arraybackslash}p{8.0cm}@{}}
\toprule
Opción & Valores & Significado \\
\midrule
\code{-{}-bmin}, \code{-{}-bmax}, \code{-{}-nbins} & malla & Malla de la
  variable distribuida (T para \code{bhf}, mm/s para \code{quad} y
  \code{delta}). \\
\code{-{}-baseline}, \code{-{}-slope} & valores & Nivel y pendiente del fondo,
  iniciales (o fijos). \\
\code{-{}-no-fit-baseline}, \code{-{}-no-fit-slope} & --- & Mantienen fijo el
  nivel o la pendiente del fondo. Valen para todas las formas. \\
\code{-{}-center}, \code{-{}-vmax} & valores & Centro de doblado y velocidad
  máxima; por defecto se toman de los ficheros \file{.RES}/\file{.JOB}/\file{.PLT}
  o se buscan por $\chi^2$ ($v_{\max} = 12$~mm/s). \\
\code{-{}-norm-percentile} & \% & Percentil con el que se normalizan las
  cuentas dobladas. \\
\code{-{}-sextet-pattern}, \code{-{}-intensity-convention},
  \code{-{}-line-asym} & & Convenios de NORMOS, como en el ajuste discreto. \\
\bottomrule
\end{tabular}
\end{center}

\subsection{Distribuciones bidimensionales}

\code{-{}-dist-2d} ajusta una distribución bidimensional (anexo~\ref{ap:dist-2d})
en una malla de \code{-{}-nbins} $\times$ \code{-{}-nbins-quad} celdas. La
pareja de variables la elige \code{-{}-pair}, como los modos 2D de la GUI:
\code{bhf-quad} para $P(\BHF,\dEQ)$ (por defecto), \code{delta-quad} para
$P(\mathrm{IS},\dEQ)$ y \code{bhf-delta} para $P(\BHF,\mathrm{IS})$. La primera
variable va entre \code{-{}-bmin} y \code{-{}-bmax}, y la segunda entre
\code{-{}-qmin} y \code{-{}-qmax}; si no se dan, se usan los rangos de la GUI
($\dEQ$ de $-1$ a $1$~mm/s con campo y de $0$ a $3$~mm/s sin él; IS de $-1$ a
$2$~mm/s). Sin $\BHF$ en la pareja, el kernel usa el campo de
\code{-{}-fixed-bhf} ($0$, dobletes, por defecto). La forma elige el método:

\begin{itemize}[nosep]
  \item \code{histograma} (por defecto): histograma regularizado. La
        regularización a lo largo de $\BHF$ es \code{-{}-alpha}, y a lo largo
        de $\dEQ$, \code{-{}-alpha-quad} (igual a \code{-{}-alpha} si no se
        da).
  \item \code{gaussiana} o \code{vbf}: una o \code{-{}-vbf-components}
        gaussianas bivariantes $(A, \mu_x, \mu_y, \sigma_x, \sigma_y, \rho)$,
        sin $\alpha$. \code{-{}-vbf-sigma-min}/\code{-{}-vbf-sigma-max} acotan
        $\sigma_x$ y \code{-{}-vbf-sigma-min-y}/\code{-{}-vbf-sigma-max-y},
        $\sigma_y$.
\end{itemize}

La forma binomial y \code{-{}-scan-alpha} no están disponibles en 2D.

\subsection{Elección de \texorpdfstring{$\alpha$}{alfa}}

\code{-{}-scan-alpha} repite el ajuste del histograma para
\code{-{}-alpha-count} valores de $\alpha$, espaciados logarítmicamente entre
\code{-{}-alpha-min} y \code{-{}-alpha-max}, con el mismo kernel y las mismas
opciones que el ajuste principal. Escribe una L-curve sencilla
(\file{*\_alpha\_scan.dat}, y un PNG con \code{-{}-plot}) y el valor sugerido de
la esquina en el resumen (\code{suggested\_alpha\_lcurve}).

\subsection{Ficheros de salida}

En los ajustes 1D se escriben tres ficheros junto al de entrada (o con el
prefijo de \code{-{}-out-prefix}): \file{*\_spectrum.dat} (velocidad, datos,
ajuste, residuo, cuentas dobladas), \file{*\_distribution.dat} (malla, $P$, $P$
normalizada) y \file{*\_summary.json} (todos los parámetros escalares y las
métricas; con VBF incluye los tríos $(A_j,\mu_j,\sigma_j)$ ajustados y las
cotas de $\sigma$ aplicadas, \code{vbf\_sigma\_bounds}). Con
\code{-{}-plot} se añade un PNG con el espectro y la distribución.

En modo 2D la distribución se escribe como matriz en
\file{*\_distribution2d.dat}, con sus dos ejes en \file{*\_<x>\_axis.dat}
y \file{*\_<y>\_axis.dat} (\file{*\_bhf\_axis.dat} y \file{*\_quad\_axis.dat}
con la pareja por defecto), y el resumen incluye la pareja y los momentos
$\langle\BHF\rangle$, $\langle\dEQ\rangle$, las anchuras y la correlación
$\BHF$--$\dEQ$; con las formas gaussianas, también las gaussianas ajustadas
(\code{vbf2d\_components}) y las cotas de $\sigma$ de cada eje
(\code{vbf2d\_sigma\_bounds}).
